\section{DFS}

\subsection{Домашнее задание}
\begin{enumerate}
  \item
    Пусть дано дерево $T = \langle V, E \rangle$. Для каждого ребра вычислить, сколько
    простых путей проходит через него. $\O(|V| + |E|)$.

  \item
    Корневое дерево $T$ на $V$ вершинах задается массивом
    из $V$ элементов. Все вершины пронумерованы. Для каждой вершины
    в массиве указан его родитель.  Для корня $r$ значение в массиве
    равно -1. Требуется определить, как будет выглядеть новое представление
    дерева, если корень $r$ сменить на корень $q$. Разрешается использовать
    $\O(1)$ дополнительной памяти. Менять массив можно. Время $\O(V)$.

  \item
    За $\O(|V| + |E|)$ найдите в \textbf{неорграфе}:
    \begin{enumerate}
      \item все \textbf{ребра}, через которые проходит любой путь $a \leadsto b$.
      \item все \textbf{вершины}, через которые проходит любой путь $a \leadsto b$.
    \end{enumerate}

  \item
    Дано дерево $T = \langle V, E \rangle$.
    \begin{enumerate}
      \item
        Найдите центроид дерева --- вершину $c$, для которой минимальна \\
        $\sum_{v \in V} dist(v, c)$. $\O(|V|)$.

      \item
        Найдите 2-центроид дерева --- пару вершин $(a, b)$, для которой минимальна \\
        $\sum_{v \in V}\min(dist(v, a), dist(v, b))$). $\O(|V|^2)$.
    \end{enumerate}

\subsection*{Дополнительные задачи}

  \item
    Для заданного ориентированного графа $G = \langle V, E \rangle$ посчитать
    минимальное число ребер, которые нужно добавить в граф, чтобы
    он стал сильно связным. Время $\O(|V|^3)$.

   \item
    Дано дерево $T = \langle V, E \rangle$. Найдите 2-центроид за $\O(|V|)$.

\end{enumerate}

\newpage

\section{Решения}

\begin{enumerate}

% Пусть дано дерево $T = \langle V, E \rangle$. Для каждого ребра вычислить, сколько
% простых путей проходит через него. $\O(|V| + |E|)$.

\item 
    
    Для решения задачи 
    

\item \dots

\item \dots

\item \dots

\subsection*{Дополнительные задачи}

\item \dots

\item \dots

\end{enumerate}
